% ============================================================
%  Abstract.tex  –  Game Tracker
% ============================================================

\chapter*{Resumen}
\addcontentsline{toc}{chapter}{Resumen}

Game Tracker es una aplicación móvil diseñada para que los usuarios gestionen
su colección personal de videojuegos y compartan reseñas dentro de una
comunidad. El sistema permite registrar juegos con atributos como nombre,
franquicia, categoría y calificación; marcar el estado de progreso de cada
título y escribir notas con soporte de respuestas anidadas. Otros usuarios
pueden reaccionar a esas notas mediante un sistema de emociones configurables.

El proyecto sigue una arquitectura cliente-servidor: el backend está construido
en Java con Spring Boot, utiliza MySQL como motor de base de datos y se despliega
en la nube a través de Render y Aiven. El frontend es una aplicación móvil
desarrollada en Flutter que se comunica con el servidor mediante una API REST
protegida con JSON Web Tokens (JWT).

A lo largo del desarrollo se aplicaron cuatro patrones de diseño reconocidos
—MVC, Singleton, Strategy y ChangeNotifier— con el objetivo de mantener el
código organizado, extensible y fácil de probar. El resultado es un sistema
funcional que cubre el ciclo completo de gestión de un catálogo de videojuegos,
desde el registro inicial hasta las interacciones de la comunidad.